\documentclass[11pt]{article}

\usepackage[T1]{fontenc}
\usepackage[margin=1in]{geometry}
\usepackage{amsmath,amssymb,amsthm}
\usepackage{booktabs}
\usepackage{array}
\usepackage{enumitem}
\usepackage{xcolor}
\usepackage{hyperref}

\emergencystretch=2em

\hypersetup{
  colorlinks=true,
  linkcolor=blue!55!black,
  citecolor=blue!55!black,
  urlcolor=blue!55!black,
  pdftitle={Non-Pseudo-Arithmetic Hyperbolic 4-Manifolds},
  pdfauthor={Alexander Kolpakov}
}

\numberwithin{equation}{section}

\newtheorem{theorem}{Theorem}[section]
\newtheorem{proposition}[theorem]{Proposition}
\newtheorem{lemma}[theorem]{Lemma}
\newtheorem{corollary}[theorem]{Corollary}
\newtheorem{fact}[theorem]{Basic fact}
\theoremstyle{definition}
\newtheorem{definition}[theorem]{Definition}
\newtheorem{example}[theorem]{Example}
\theoremstyle{remark}
\newtheorem{remark}[theorem]{Remark}

\newcommand{\Q}{\mathbb{Q}}
\newcommand{\R}{\mathbb{R}}
\newcommand{\C}{\mathbb{C}}
\newcommand{\HH}{\mathbb{H}}
\newcommand{\calG}{\mathcal{G}}
\newcommand{\calP}{\mathcal{P}}
\newcommand{\Gal}{\operatorname{Gal}}
\newcommand{\tr}{\operatorname{tr}}
\newcommand{\Ad}{\operatorname{Ad}}
\newcommand{\Isom}{\operatorname{Isom}}
\newcommand{\PO}{\operatorname{PO}}
\newcommand{\diag}{\operatorname{diag}}
\newcommand{\minpoly}{\operatorname{minpoly}}
\newcommand{\Spec}{\operatorname{Spec}}
\newcommand{\vol}{\operatorname{vol}}
\newcommand{\code}[1]{\texttt{#1}}

\title{Non-Pseudo-Arithmetic Hyperbolic \(4\)-Manifolds}
\author{Alexander Kolpakov}
\date{24 September 2026}

\begin{document}

\maketitle

\begin{abstract}
We give three explicit finite-volume hyperbolic Coxeter lattices in dimension
four whose adjoint trace fields are not totally real.  They are the reflection
groups of the Ma--Zheng polytopes
\(P_{11,6},P_{11,7},P_{11,8}\).  The obstruction is contained in two cyclic
products of the doubled unit-normal Gram matrix \(A=2G\).  Those products
recover elements \(s=\sqrt5\) and \(u\) satisfying
\[
 (s-1)u^2+2(s-3)u+s-7=0.
\]
After the embedding \(s\mapsto-\sqrt5\), this quadratic has discriminant
\(8(1-\sqrt5)<0\).  The cyclic-product field therefore has a complex place.
Vinberg's identification of that field with the adjoint trace field and the
Emery--Mila total-reality obstruction imply that the three lattices are not
pseudo-arithmetic.  Torsion-free orientation-preserving finite-index
subgroups yield orientable noncompact finite-volume hyperbolic
\(4\)-manifolds with the same obstruction.

The exposition includes all specialized definitions used in the argument,
an exact source-to-matrix transcription, a proof of every new algebraic
step, independent exact certificates, and an explicit register of imported
theorems.  A pinned Lean~4 development kernel-checks the radical identity,
the number-field obstruction, the main implication, and the finite-index
corollary.  Statements imported from geometry and the literature are
represented in Lean as local hypotheses rather than asserted as axioms or
reproved here.
\end{abstract}

\tableofcontents

\section{Introduction and complete reader's primer}
\label{sec:introduction}

\subsection{The question and the answer}

A \emph{lattice} in the isometry group of real hyperbolic \(n\)-space is a
discrete group whose quotient has finite hyperbolic volume.  Arithmetic
lattices are constructed from algebraic groups over number fields.  Many
important nonarithmetic lattices are assembled by gluing arithmetic pieces.
Emery and Mila introduced \emph{pseudo-arithmeticity} to give one common
algebraic framework for arithmetic lattices, quasi-arithmetic lattices, and
the familiar gluing constructions \cite{EmeryMila}.  They asked whether every
lattice in \(\PO(n,1)\), for \(n>3\), is pseudo-arithmetic and, more weakly,
whether every such lattice has totally real adjoint trace field
\cite[Questions~1.8 and~1.10]{EmeryMila}.

Both questions have negative answers in dimension four.

\begin{theorem}[Main theorem]
\label{thm:main}
For \(j\in\{6,7,8\}\), let \(P_{11,j}\) be the finite-volume hyperbolic
Coxeter \(4\)-polytope in the Ma--Zheng census recorded in
Table~\ref{tab:three-polytopes}, and let
\(\Gamma_{11,j}<\PO(4,1)\) be its reflection group.  Then the adjoint trace
field of \(\Gamma_{11,j}\) is not totally real.  Consequently
\(\Gamma_{11,j}\) is not pseudo-arithmetic.
\end{theorem}

\begin{corollary}[Manifold form]
\label{cor:manifolds}
There exist orientable, noncompact, finite-volume hyperbolic
\(4\)-manifolds whose fundamental groups are not pseudo-arithmetic.
\end{corollary}

Ma and Zheng classify all \(331\) finite-volume hyperbolic Coxeter
\(4\)-polytopes with seven facets \cite{MaZheng}.  The present examples are
three consecutive entries in that census.  The proof needs only four entries
of their Gram matrices and one common radical.  It does not require a
numerical approximation, a primitive element for the entire field, or a
classification of all embeddings of that field.

Emery--Mila checked the higher-dimensional Coxeter examples then available
to them.  The final 2026 version of
Belolipetsky--Bogachev--Kolpakov--Slavich still described the known examples
in \(\PO(n,1)\), \(n>3\), as pseudo-arithmetic \cite{BBKS}.  A targeted
literature audit through 24 September 2026 found no earlier recorded
non-pseudo-arithmetic lattice in these dimensions.  We therefore use the
qualified statement that we know of no previous example; a bounded search is
not a proof of absolute priority.

\subsection{All specialized language used in the proof}
\label{subsec:primer}

This subsection is a compact primer.  Later sections repeat the definitions
at the point where calculations use them.

\paragraph{Hyperbolic space and lattices.}
On \(\R^{n+1}\), put
\[
 \langle x,y\rangle_{n,1}
   =x_1y_1+\cdots+x_ny_n-x_{n+1}y_{n+1}.
\]
The upper sheet
\[
 \HH^n=\{x\in\R^{n+1}:\langle x,x\rangle_{n,1}=-1,
                       \ x_{n+1}>0\}
\]
is the hyperboloid model of real hyperbolic space.  The group preserving the
form acts by isometries; after quotienting by its scalar center one obtains
the projective orthogonal group \(\PO(n,1)\).  A subgroup
\(\Gamma<\Isom(\HH^n)\) is \emph{discrete} if the identity is isolated in
\(\Gamma\), and it is a \emph{lattice} if
\(\Gamma\backslash\HH^n\) has finite volume.  The quotient is an orbifold in
general and a manifold when \(\Gamma\) is torsion-free.

\paragraph{Coxeter polytopes and Gram matrices.}
A convex hyperbolic polytope is \emph{Coxeter} when every dihedral angle is
\(\pi/m\) for an integer \(m\ge2\).  Reflections in its facets generate a
discrete group for which the polytope is a fundamental domain.  If the
polytope has finite volume, this reflection group is a lattice.  Choose an
outward unit spacelike normal \(e_i\) to each facet.  The matrix
\(G=(\langle e_i,e_j\rangle_{n,1})\) is the \emph{unit-normal Gram matrix}.
Its diagonal entries are one.  Off the diagonal one has
\[
 g_{ij}=
 \begin{cases}
 -\cos(\pi/m_{ij}),&\text{the facets meet at angle }\pi/m_{ij},\\
 -1,&\text{the facets are parallel},\\
 -\cosh(\rho_{ij}),&\text{they are disjoint at distance }\rho_{ij}.
 \end{cases}
\]
We use \(A=2G\).  A \emph{cyclic product} is a product
\(A_{i_0i_1}A_{i_1i_2}\cdots A_{i_mi_0}\).  The field generated over \(\Q\)
by every such product is the \emph{Vinberg field} \(k(P)\).

\paragraph{Number fields and total reality.}
A \emph{number field} is a finite extension of \(\Q\).  An embedding of a
number field \(K\) is an injective field homomorphism \(K\hookrightarrow\C\)
fixing \(\Q\).  The field \(K\) is \emph{totally real} when every such
embedding has image in \(\R\).  Thus one nonreal embedding is a complete
certificate that \(K\) is not totally real.  We use the standard extension
theorem: if \(L\subset K\) are number fields, every embedding
\(L\hookrightarrow\C\) extends to an embedding \(K\hookrightarrow\C\).
This follows from the existence of a complex root of the transported minimal
polynomial; a proof is recalled in Proposition~\ref{prop:embedding-extension}.

\paragraph{Adjoint trace fields.}
For a matrix Lie group, conjugation differentiates to a linear action
\(\Ad\) on its Lie algebra.  The \emph{adjoint trace field} of a lattice
\(\Gamma\) is
\[
 k(\Gamma)=
 \Q\bigl(\{\tr(\Ad\gamma):\gamma\in\Gamma\}\bigr)\subset\R.
\]
It is unchanged on passing to a finite-index subgroup.  Vinberg's
field-of-definition theorem identifies this field, for a finite-volume
Coxeter reflection group, with the cyclic-product field \(k(P)\).  This
identification is an imported theorem, not a new result proved here.

\paragraph{Pseudo-arithmeticity.}
Vinberg associates to a Zariski-dense lattice \(\Gamma\) an algebraic group
over \(k(\Gamma)\), called its ambient group.  An algebraic group over a
totally real field \(k\) is \emph{admissible} for hyperbolic \(n\)-space when
one real factor is the required noncompact projective orthogonal group and
the other real factors are compact.  In Emery--Mila's definition, a lattice
is \emph{pseudo-arithmetic} when its ambient group is obtained from such an
admissible group by scalar extension through a finite real multiquadratic
extension \(K/k\).  Their definition implies that \(K=k(\Gamma)\) is totally
real.  This last necessary condition is the only part of
pseudo-arithmeticity used in the proof.

\paragraph{Formal verification terminology.}
Lean checks a theorem only from its declarations, imported library theorems,
and explicitly supplied hypotheses.  In this release, the new algebraic
argument is proved inside Lean.  A literature result such as Vinberg's theorem
is represented by a named hypothesis with the same logical type needed by
the manuscript.  Lean checks every implication after that boundary, but this
does not constitute a new formal proof of Vinberg's theorem.  No
project-defined \code{axiom}, \code{opaque}, \code{sorry}, or \code{admit}
is used.

\subsection{The certificate in one paragraph}
\label{subsec:certificate-preview}

Let \(G\) be any of the three unit-normal Gram matrices and put \(A=2G\).
With zero-based facet indices, the common entries
\[
 A_{12}=-\frac{1+\sqrt5}{2},\qquad
 A_{03}=-1,\qquad A_{35}=-\sqrt2,\qquad A_{50}=-2a
\]
give two cyclic products
\[
 c_{12}=A_{12}A_{21}=\frac{3+\sqrt5}{2},\qquad
 c_{035}=A_{03}A_{35}A_{50}=-2\sqrt2\,a.
\]
Hence the Vinberg field contains
\(s=2c_{12}-3=\sqrt5\) and
\(u=-c_{035}/2=\sqrt2a\).  The exact radical for \(a\) implies
\[
 (s-1)u^2+2(s-3)u+s-7=0.
\]
At the conjugate embedding of \(\Q(\sqrt5)\), this quadratic has negative
discriminant.  It therefore produces a nonreal embedding of
\(\Q(s,u)\), which extends to the whole Vinberg field.  This is the entire
new obstruction.

\subsection{Logical status of the ingredients}
\label{subsec:logical-status}

For clarity, the proof has four layers.

\begin{enumerate}[label=\textbf{L\arabic*.},leftmargin=3em]
\item \emph{Elementary proof:} the matrix entries, two products, radical
  elimination, discriminant, and complex-place deduction are written out in
  Sections~\ref{sec:coxeter-data}--\ref{sec:main-proof}.
\item \emph{Exact computation:} independent Python and PARI/GP programs check
  the source transcription, rank, inertia, field signatures, and finite-volume
  combinatorics.  They support the input data but are not needed to trust the
  two-cycle algebra once the displayed data are accepted.
\item \emph{Lean kernel checking:} the internal implication from the radical
  formula and cycle certificate to non-total reality and non-pseudo-
  arithmeticity is formalized.  The finite-index corollary is checked from
  explicitly named geometric hypotheses.
\item \emph{External mathematics:} the census classification, Vinberg's
  field identification, the Emery--Mila necessary condition, Selberg's lemma,
  and finite-index invariance are cited rather than reproved.  The exact list
  appears in Section~\ref{sec:external-register}.
\end{enumerate}

This separation prevents an executable check of a matrix identity from being
misdescribed as a formal verification of a classification theorem.

\section{Basic field and hyperbolic facts}
\label{sec:basics}

\subsection{Embeddings, minimal polynomials, and complex places}

We record the elementary field facts at the level needed later.

\begin{definition}
Let \(K\) be a number field and \(\alpha\in K\).  The \emph{minimal
polynomial} \(\minpoly_\Q(\alpha)\) is the unique monic polynomial of least
positive degree in \(\Q[X]\) that vanishes at \(\alpha\).
\end{definition}

The uniqueness follows from Euclidean division: if two monic polynomials of
least degree vanish, their difference has smaller degree and also vanishes.
The kernel of the evaluation homomorphism
\(\Q[X]\to K\), \(p\mapsto p(\alpha)\), is therefore generated by this
polynomial.

\begin{lemma}
\label{lem:sqrt5-minpoly}
The number \(\sqrt5\) is irrational and
\[
 \minpoly_\Q(\sqrt5)=X^2-5.
\]
Its two embeddings over \(\Q\) send \(\sqrt5\) to \(\sqrt5\) and
\(-\sqrt5\), respectively.
\end{lemma}

\begin{proof}
Suppose \(\sqrt5=p/q\) for coprime positive integers \(p,q\).  Then
\(p^2=5q^2\), so the prime \(5\) divides \(p^2\), hence divides \(p\).
Write \(p=5r\).  Substitution gives \(q^2=5r^2\), so \(5\mid q\), contrary
to coprimality.  Thus \(\sqrt5\notin\Q\).  It is a root of the monic
quadratic \(X^2-5\); irrationality rules out a linear minimal polynomial,
so this quadratic is minimal.  Any embedding preserves addition,
multiplication, and rational numbers, and hence must send \(\sqrt5\) to a
root of \(X^2-5\).  The two roots are \(\pm\sqrt5\), and both choices define
an embedding of \(\Q(\sqrt5)\).
\end{proof}

\begin{proposition}[Extension of embeddings]
\label{prop:embedding-extension}
Let \(L\subset K\) be number fields and let
\(\sigma:L\hookrightarrow\C\) be an embedding.  Then \(\sigma\) extends to
an embedding \(\widetilde\sigma:K\hookrightarrow\C\).
\end{proposition}

\begin{proof}
This is the standard embedding-extension theorem; we recall its mechanism.
Because number fields have characteristic zero, \(K/L\) is finite and
separable.  By the primitive-element theorem there is \(\theta\in K\) with
\(K=L(\theta)\).  Let \(p\in L[X]\) be the minimal polynomial of \(\theta\).
Apply \(\sigma\) coefficientwise to obtain \(p^\sigma\in\C[X]\).  The
fundamental theorem of algebra supplies a root \(z\in\C\) of
\(p^\sigma\).  The map
\[
 L[X]\longrightarrow\C,
 \qquad \sum_i a_iX^i\longmapsto\sum_i\sigma(a_i)z^i
\]
has \((p)\) in its kernel and therefore descends to a homomorphism
\(L[X]/(p)\to\C\).  Identifying \(L[X]/(p)\) with \(L(\theta)=K\) gives the
extension.  A unital homomorphism from a field has zero kernel, so it is an
embedding.
\end{proof}

\begin{lemma}[Quadratic real-root test]
\label{lem:quadratic-test}
Let \(a,b,c\in\R\), with \(a\ne0\), and put
\(D=b^2-4ac\).  The equation \(ax^2+bx+c=0\) has a real solution if and
only if \(D\ge0\).
\end{lemma}

\begin{proof}
Multiplying the equation by \(4a\) and completing the square gives
\[
 (2ax+b)^2=b^2-4ac=D.
\]
A real square is nonnegative, proving necessity.  If \(D\ge0\), the two
values \((-b\pm\sqrt D)/(2a)\) are real and substitution shows that they are
solutions.
\end{proof}

\begin{proposition}[A nonreal subfield embedding persists]
\label{prop:subfield-persists}
Let \(L\subset K\) be number fields.  If \(L\) has a nonreal embedding, then
\(K\) is not totally real.
\end{proposition}

\begin{proof}
Let \(\sigma:L\hookrightarrow\C\) have image not contained in \(\R\).
Proposition~\ref{prop:embedding-extension} extends it to
\(\widetilde\sigma:K\hookrightarrow\C\).  If the image of
\(\widetilde\sigma\) were contained in \(\R\), then so would be the image
of its restriction \(\sigma\), a contradiction.
\end{proof}

\subsection{The hyperboloid model and reflecting hyperplanes}

Let \(V=\R^{n+1}\) with the form \(\langle\ ,\ \rangle_{n,1}\) defined in
Subsection~\ref{subsec:primer}.  A vector is spacelike, null, or timelike
according as its norm is positive, zero, or negative.  If \(e\) is a unit
spacelike vector, then
\[
 H_e=\{x\in\HH^n:\langle x,e\rangle_{n,1}=0\}
\]
is a totally geodesic hyperplane, and
\[
 r_e(x)=x-2\langle x,e\rangle_{n,1}e
\]
is the reflection fixing \(H_e\) pointwise.  Indeed,
\(r_e(e)=-e\), \(r_e\) fixes the orthogonal complement of \(e\), and direct
expansion shows
\(\langle r_e(x),r_e(y)\rangle_{n,1}=\langle x,y\rangle_{n,1}\).

For a polytope described as the intersection of half-spaces
\(\langle x,e_i\rangle_{n,1}\le0\), the vectors \(e_i\) are outward normals.
The inner products of pairs of normals give the Gram formulas in
Subsection~\ref{subsec:primer}; these are the standard two-dimensional
Lorentzian calculations for intersecting, asymptotic, and ultraparallel
hyperplanes \cite[Chapter~3]{Ratcliffe}.

\begin{definition}
A Coxeter matrix has diagonal entries one and off-diagonal labels
\(m_{ij}\in\{2,3,\ldots\}\cup\{\infty\}\).  Its Coxeter group is generated
by involutions \(r_i\) with relations \((r_ir_j)^{m_{ij}}=1\) when
\(m_{ij}<\infty\).  For a hyperbolic Coxeter polytope, these generators are
the geometric facet reflections.
\end{definition}

The fact that the geometric reflection group is discrete and that the
polytope is a fundamental domain is the Poincar\'e--Coxeter theorem.  We use
it as standard background.  Finite volume of the fundamental domain is
exactly finite covolume of the group.

\subsection{From an orbifold to a manifold}
\label{subsec:orbifold-manifold}

Two elementary finite-index observations will be used.  First, the kernel of
the orientation character \(\Gamma\to\{\pm1\}\) has index at most two.
For a group generated by hyperplane reflections, each generator reverses
orientation, so the kernel has index exactly two.  Second, a subgroup of
finite index \(d\) gives a degree-\(d\) orbifold cover and multiplies volume
by \(d\).  Thus finite volume is preserved in both directions.

The following external theorem removes torsion.

\begin{theorem}[Selberg's lemma]
\label{thm:selberg}
Every finitely generated subgroup of \(\mathrm{GL}_N(F)\), where \(F\) has
characteristic zero, has a torsion-free subgroup of finite index.
\end{theorem}

We cite Selberg's original result \cite{Selberg}.  A Coxeter group with
finitely many facets is finitely generated, and its reflection
representation is linear over a characteristic-zero field, so the theorem
applies.  A torsion-free discrete group acts freely on \(\HH^n\): an element
fixing a point lies in a compact point stabilizer, and discreteness makes the
cyclic subgroup finite, hence trivial.  The quotient is therefore a
manifold.  Passing inside the orientation-preserving subgroup makes it
orientable.  Finally, a finite cover of a noncompact finite-volume quotient
is noncompact: if the cover were compact, its continuous image would be
compact, contrary to noncompactness of the base.

\section{Adjoint trace fields and pseudo-arithmeticity}
\label{sec:trace-fields}

\subsection{The adjoint representation and ambient group}

Let \(\mathbf G\) be a real algebraic matrix group with Lie algebra
\(\mathfrak g\).  Conjugation by \(g\in\mathbf G(\R)\) maps
\(X\in\mathfrak g\) to \(gXg^{-1}\).  This linear map is \(\Ad(g)\).  Its
trace is independent of a choice of basis for \(\mathfrak g\).  For a
Zariski-dense subgroup \(\Gamma\), define
\[
 k(\Gamma)=\Q\bigl(\tr(\Ad\gamma):\gamma\in\Gamma\bigr).
\]
For lattices in a connected noncompact simple Lie group, Borel density gives
the required Zariski density.  Vinberg's theory produces a form of the
ambient algebraic group over \(k(\Gamma)\), unique in the appropriate
isomorphism or isogeny class \cite{VinbergRings}.  We do not need its
construction; only the field is used below.

We also use the standard finite-index invariance
\begin{equation}
 [\Gamma:\Lambda]<\infty
 \quad\Longrightarrow\quad k(\Lambda)=k(\Gamma).
\label{eq:finite-index-field}
\end{equation}
This is part of Vinberg's commensurability theory.  One inclusion is
immediate from \(\Lambda\subset\Gamma\); equality is the nontrivial external
statement and is listed as Input~E6 in
Section~\ref{sec:external-register}.

\subsection{Pseudo-admissibility}

An algebraic \(k\)-group \(\mathbf H\) is \emph{admissible for
\(\PO(n,1)\)} when \(k\) is totally real and, among the groups obtained from
the real embeddings of \(k\), exactly one has the required noncompact
factor \(\PO(n,1)^\circ\), while the others are compact.  This is the usual
arithmetic admissibility condition.

\begin{definition}[Emery--Mila]
Let \(n>3\).  A group over a number field \(K\) is
\emph{pseudo-admissible over \(K/k\)} if it is the scalar extension
\(\mathbf H_K\) of an admissible \(k\)-group through a finite, possibly
trivial, real multiquadratic extension
\[
 K=k(\sqrt{d_1},\ldots,\sqrt{d_t}).
\]
A lattice \(\Gamma<\PO(n,1)\) is \emph{pseudo-arithmetic} when its ambient
group is pseudo-admissible \cite[Definitions~1.6--1.7]{EmeryMila}.
\end{definition}

The trivial case \(t=0\) includes quasi-arithmetic lattices.  The adjective
``real'' in this definition is the archimedean condition ensuring that the
extension has only real embeddings over the relevant real places.

\begin{proposition}[Total-reality obstruction]
\label{prop:total-real-obstruction}
If \(n>3\) and \(\Gamma<\PO(n,1)\) is pseudo-arithmetic, then its adjoint
trace field \(k(\Gamma)\) is totally real.
\end{proposition}

\begin{proof}
This is the immediate necessary condition recorded by Emery--Mila
\cite[Section~1.4]{EmeryMila}.  In their definition the base field is totally
real by admissibility and \(K/k\) is a real multiquadratic extension.  Thus
every embedding \(K\hookrightarrow\C\) has image in \(\R\).  Since
\(K=k(\Gamma)\), the adjoint trace field is totally real.
\end{proof}

The contrapositive is what we use:
\begin{equation}
 k(\Gamma)\text{ not totally real}
 \quad\Longrightarrow\quad
 \Gamma\text{ not pseudo-arithmetic}.
\label{eq:total-real-contrapositive}
\end{equation}

\subsection{Vinberg's cyclic-product field}
\label{subsec:vinberg}

Let \(P\subset\HH^n\) be a finite-volume Coxeter polytope with unit-normal
Gram matrix \(G=(g_{ij})\), and put \(A=2G\).  For a closed sequence of facet
indices \(i_0,i_1,\ldots,i_m,i_0\), its cyclic product is
\[
 A_{i_0i_1}A_{i_1i_2}\cdots A_{i_{m-1}i_m}A_{i_mi_0}.
\]
Repeated indices are allowed, and reversing the cycle gives the same value
because \(A\) is symmetric.  Define
\[
 k(P)=\Q(\text{all cyclic products in }A).
\]

\begin{theorem}[Vinberg]
\label{thm:vinberg-cycles}
If \(P\) has finite volume and \(\Gamma(P)\) is its reflection group, then
\[
 k(\Gamma(P))=k(P).
\]
\end{theorem}

This is the cyclic-product form of Vinberg's field-of-definition theorem
\cite{VinbergRings}; see also \cite[Section~4.1]{EmeryMila} and
\cite[Section~3.1]{GuoMaZangZheng}.  It is an external input.  The factor two
is not cosmetic: Vinberg's convention uses the bilinear Gram matrix with
diagonal two.  Starting from unit normals, that matrix is \(A=2G\).

\begin{remark}[Entry field versus trace field]
The field generated by all individual entries of \(G\) can be strictly
larger than \(k(P)\).  An individual entry changes if a normal vector is
rescaled, whereas closed products cancel these rescalings.  Accordingly,
showing that a radical occurs in one Gram entry does not show that it lies in
the adjoint trace field.  Every field inclusion below is obtained from an
explicit closed product.
\end{remark}

\section{The three Coxeter polytopes and their exact matrices}
\label{sec:coxeter-data}

\subsection{How the census vector is read}

We number the seven facets \(0,1,\ldots,6\).  The \(21\) unordered pairs are
listed in lexicographic blocks:
\begin{equation}
\begin{split}
 &(01,02,03,04,05,06)\mid(12,13,14,15,16)\mid(23,24,25,26)\\
 &\hspace{43mm}\mid(34,35,36)\mid(45,46)\mid(56).
\end{split}
\label{eq:pair-order}
\end{equation}
An integer \(m\ge2\) means dihedral angle \(\pi/m\), the symbol \(0\)
means parallel facets, and a letter means an ultraparallel pair.  Thus the
initial block \(2232ab\) means
\[
 m_{01}=2,\quad m_{02}=2,\quad m_{03}=3,\quad m_{04}=2,
 \quad g_{05}=-a,\quad g_{06}=-b.
\]

The relevant published rows are as follows.  The displayed volume is the
volume of the Coxeter fundamental polytope.

\begin{table}[ht]
\centering
\small
\begin{tabular}{@{}c l c c@{}}
\toprule
polytope & Coxeter vector & cusps & volume \\
\midrule
\(P_{11,6}\) & \(\mathtt{2232ab\ 52322\ 2220\ 442\ 22\ 2}\)
    & \(1\) & \(13\pi^2/864\) \\
\(P_{11,7}\) & \(\mathtt{2232ab\ 52323\ 2220\ 442\ 22\ 2}\)
    & \(1\) & \(\pi^2/54\) \\
\(P_{11,8}\) & \(\mathtt{2232ab\ 52324\ 2220\ 442\ 22\ 2}\)
    & \(2\) & \(19\pi^2/864\) \\
\bottomrule
\end{tabular}
\caption{The three Ma--Zheng census rows used in the proof.}
\label{tab:three-polytopes}
\end{table}

The assertion that these rows define finite-volume polytopes is external
Input~E1.  Section~\ref{sec:certification} describes independent exact checks,
including a direct finite-volume check for \(P_{11,8}\).

\subsection{Radicals and a uniform matrix}

Put
\[
 s=\sqrt5,\qquad r=\sqrt2,
\]
and let \(q_6=2,q_7=3,q_8=4\).  The common first dotted entry is
\begin{equation}
 a=\frac12\sqrt{7+s+2\sqrt{2+2s}}.
\label{eq:a-radical}
\end{equation}
The second dotted entries are
\begin{equation}
\begin{aligned}
 b_6&=\sqrt{\frac92+2s+\sqrt{38+17s}},\\
 b_7&=r+\frac{\sqrt{10}}2+\frac12\sqrt{29+13s},\\
 b_8&=\frac14\left(2+3r+2s+rs
       +2\sqrt{19+14r+9s+6rs}\right).
\end{aligned}
\label{eq:b-radicals}
\end{equation}
All square roots mean the positive real root.  Their radicands are positive,
and the resulting \(a,b_j\) exceed one, as required for ultraparallel
facets.  Only \(a\) is used in the obstruction.

For reference, the complete unit-normal Gram matrix is
\begin{equation}
G_j=
\begin{pmatrix}
1&0&0&-\tfrac12&0&-a&-b_j\\
0&1&-\tfrac{1+s}{4}&0&-\tfrac12&0&-\cos(\pi/q_j)\\
0&-\tfrac{1+s}{4}&1&0&0&0&-1\\
-\tfrac12&0&0&1&-\tfrac r2&-\tfrac r2&0\\
0&-\tfrac12&0&-\tfrac r2&1&0&0\\
-a&0&0&-\tfrac r2&0&1&0\\
-b_j&-\cos(\pi/q_j)&-1&0&0&0&1
\end{pmatrix}.
\label{eq:full-gram}
\end{equation}
To check the transcription, read each row against \eqref{eq:pair-order}.
For example, \(m_{12}=5\) gives
\(g_{12}=-\cos(\pi/5)=-(1+s)/4\); \(m_{26}=0\) gives \(g_{26}=-1\);
and \(m_{34}=m_{35}=4\) gives \(-r/2\).  The elementary special-angle
identities used here are
\[
 \cos(\pi/3)=\frac12,\qquad
 \cos(\pi/4)=\frac{\sqrt2}{2},\qquad
 \cos(\pi/5)=\frac{1+\sqrt5}{4}.
\]
The first two follow from an equilateral triangle and an isosceles right
triangle.  For the third, put \(z=e^{i\pi/5}\) and
\(x=z+z^{-1}=2\cos(\pi/5)\).  Since \(z^5=-1\) and \(z\ne-1\), division of
\(z^5+1\) by \(z+1\) gives
\(z^4-z^3+z^2-z+1=0\).  After division by \(z^2\), this becomes
\[
 (z^2+z^{-2})-(z+z^{-1})+1=0.
\]
Because \(z^2+z^{-2}=x^2-2\), we get \(x^2-x-1=0\).  Positivity selects
\(x=(1+\sqrt5)/2\).

\subsection{The two cyclic products}
\label{subsec:two-cycles}

Put \(A_j=2G_j\).  Four common entries of \(A_j\) are
\begin{equation}
 A_{12}=-\frac{1+s}{2},\qquad
 A_{03}=-1,\qquad A_{35}=-r,
 \qquad A_{50}=-2a.
\label{eq:four-entries}
\end{equation}
The first closed path has two directed edges, \(1\to2\to1\), and gives
\begin{equation}
\begin{aligned}
 c_{12}
 &=A_{12}A_{21}
   =\left(-\frac{1+s}{2}\right)^2\\
 &=\frac{1+2s+s^2}{4}
   =\frac{6+2s}{4}
   =\frac{3+s}{2},
\end{aligned}
\label{eq:cycle12}
\end{equation}
where \(s^2=5\).  Thus
\begin{equation}
 s=2c_{12}-3\in k(P_{11,j}).
\label{eq:recover-s}
\end{equation}
The second closed path is \(0\to3\to5\to0\):
\begin{equation}
 c_{035}=A_{03}A_{35}A_{50}=(-1)(-r)(-2a)=-2ra.
\label{eq:cycle035}
\end{equation}
Consequently
\begin{equation}
 u:=ra=-\frac12c_{035}\in k(P_{11,j}).
\label{eq:recover-u}
\end{equation}
These equations prove membership in the cyclic-product field directly.  No
individual-entry argument is being used.

\section{The radical identity and the complex place}
\label{sec:field-obstruction}

\subsection{Deriving the quadratic relation}

We now prove the one algebraic identity behind all three cases.

\begin{lemma}[Exact radical elimination]
\label{lem:radical-relation}
Let \(s=\sqrt5\), let \(a\) be given by \eqref{eq:a-radical}, and let
\(u=\sqrt2a\).  Then
\begin{equation}
 (s-1)u^2+2(s-3)u+s-7=0.
\label{eq:cycle-relation}
\end{equation}
\end{lemma}

\begin{proof}
Set
\[
 t=\sqrt{2+2s}>0.
\]
By the definitions of \(a\) and \(u\),
\begin{equation}
 u^2=2a^2
 =2\cdot\frac14(7+s+2t)
 =\frac{7+s}{2}+t.
\label{eq:u-square}
\end{equation}
We claim first that
\begin{equation}
 (s-1)u=3-s+t.
\label{eq:linear-radical}
\end{equation}
Both sides are positive: \(s>1\) and \(u>0\) for the left side, while
\(s<3\) and \(t>0\) for the right side.  It is therefore enough to prove
that their squares are equal.  Using \(s^2=5\), \(t^2=2+2s\), and
\eqref{eq:u-square}, expansion gives
\[
\begin{aligned}
 ((s-1)u)^2
 &=(s-1)^2\left(\frac{7+s}{2}+t\right)\\
 &=(6-2s)\left(\frac{7+s}{2}+t\right)\\
 &=16-4s+(6-2s)t,
\end{aligned}
\]
because
\((6-2s)(7+s)/2=(42-8s-2s^2)/2=16-4s\).  On the other hand,
\[
\begin{aligned}
 (3-s+t)^2
 &=(3-s)^2+2(3-s)t+t^2\\
 &=14-6s+2(3-s)t+2+2s\\
 &=16-4s+(6-2s)t.
\end{aligned}
\]
The positive quantities have equal squares, proving
\eqref{eq:linear-radical}.

Rearrange it as
\begin{equation}
 t=(s-1)u+s-3.
\label{eq:t-linear}
\end{equation}
Squaring, substituting \(t^2=2+2s\), and moving everything to the left gives
\[
 ((s-1)u+s-3)^2-(2+2s)=0.
\]
A direct expansion factors the left side as
\[
 (s-1)\bigl((s-1)u^2+2(s-3)u+s-7\bigr).
\]
Because \(s\ne1\), division by \(s-1\) proves the claimed relation.
\end{proof}

\begin{remark}
The expanded proof uses only positivity of the named square roots,
\(s^2=5\), and ring arithmetic.  The Lean theorem
\begin{center}
\code{explicit\_radical\_cycle\_relation}
\end{center}
checks the same identity from the nested square-root definitions.  The
intermediate theorem
\begin{center}
\code{cycle\_relation\_of\_radical\_identities}
\end{center}
checks the final factorization over an arbitrary characteristic-zero field.
\end{remark}

\subsection{Negative discriminant at the conjugate place}

\begin{lemma}[Non-total-reality certificate]
\label{lem:non-real-subfield}
Let \(s=\sqrt5\) and let \(u\) satisfy \eqref{eq:cycle-relation}.  Then
\(L=\Q(s,u)\) is not totally real.
\end{lemma}

\begin{proof}
View \eqref{eq:cycle-relation} as \(f(u)=0\), where
\[
 f(X)=(s-1)X^2+2(s-3)X+s-7\in\Q(s)[X].
\]
Its leading coefficient is nonzero because \(s\ne1\).  Its discriminant is
\begin{equation}
\begin{aligned}
 \Delta
 &=\bigl(2(s-3)\bigr)^2-4(s-1)(s-7)\\
 &=4\bigl((s-3)^2-(s-1)(s-7)\bigr)\\
 &=4\bigl((s^2-6s+9)-(s^2-8s+7)\bigr)\\
 &=8(1+s).
\end{aligned}
\label{eq:discriminant}
\end{equation}
Let \(\sigma:\Q(s)\hookrightarrow\R\) be the nonidentity embedding, so
\(\sigma(s)=-s\).  Applying \(\sigma\) to the coefficients gives a real
quadratic \(f^\sigma\) with discriminant
\begin{equation}
 \sigma(\Delta)=8(1-s)<0,
\label{eq:negative-discriminant}
\end{equation}
because \(s=\sqrt5>1\).  By Lemma~\ref{lem:quadratic-test},
\(f^\sigma\) has no real root.

If \(u\in\Q(s)\), then applying \(\sigma\) to \(f(u)=0\) would make the
real number \(\sigma(u)\) a root of \(f^\sigma\), a contradiction.  Hence
\(u\notin\Q(s)\).  Since \(f\) is quadratic and has \(u\) as a root, it is
the minimal polynomial of \(u\) over \(\Q(s)\).  Extend \(\sigma\) from
\(\Q(s)\) to \(L=\Q(s,u)\).  The image of \(u\) must be one of the two
roots of \(f^\sigma\); neither root is real.  Thus this extension is a
nonreal embedding of \(L\), and \(L\) is not totally real.
\end{proof}

There is a shorter intrinsic formulation.  If a number field \(K\) contains
\(s,u\) with
\[
 \minpoly_\Q(s)=X^2-5,
 \qquad (s-1)u^2+2(s-3)u+s-7=0,
\]
then the root \(-\sqrt5\) of the minimal polynomial is realized as the image
of \(s\) under some embedding \(K\hookrightarrow\C\).  If that embedding
were real, the image of \(u\) would solve the impossible conjugate quadratic.
This is exactly the formulation kernel-checked in
\code{not\_totallyReal\_of\_cycle\_minpoly}.

\section{Proofs of the main theorem and manifold corollary}
\label{sec:main-proof}

\begin{proof}[Proof of Theorem~\ref{thm:main}]
Fix \(j\in\{6,7,8\}\).  Let
\(\Gamma_j=\Gamma(P_{11,j})\) be the group generated by reflections in the
facets of \(P_{11,j}\).  By the Ma--Zheng census statement, \(P_{11,j}\) has
finite volume.  Therefore \(\Gamma_j\) is a lattice.

Let \(K_j=k(\Gamma_j)\) be its adjoint trace field.  By
Theorem~\ref{thm:vinberg-cycles},
\begin{equation}
 K_j=k(P_{11,j}).
\label{eq:vinberg-equality}
\end{equation}
The products in \eqref{eq:cycle12} and \eqref{eq:cycle035} are cyclic
products of \(A_j=2G_j\).  Equations \eqref{eq:recover-s} and
\eqref{eq:recover-u} therefore show that
\[
 s=\sqrt5\in K_j,
 \qquad u=\sqrt2a\in K_j.
\]
Consequently \(L=\Q(s,u)\) is a subfield of \(K_j\).

Lemma~\ref{lem:radical-relation} gives the quadratic relation for \(s,u\),
and Lemma~\ref{lem:non-real-subfield} gives a nonreal embedding of \(L\).
By Proposition~\ref{prop:subfield-persists}, that embedding extends to a
nonreal embedding of \(K_j\).  Thus \(K_j\) is not totally real.

If \(\Gamma_j\) were pseudo-arithmetic,
Proposition~\ref{prop:total-real-obstruction} would make \(K_j\) totally
real.  This contradicts the preceding paragraph.  Hence \(\Gamma_j\) is not
pseudo-arithmetic.  The computation was independent of \(j\), so it applies
to all three rows.
\end{proof}

\begin{proof}[Proof of Corollary~\ref{cor:manifolds}]
Choose one \(j\) and write \(\Gamma=\Gamma_{11,j}\).  Let
\(\Gamma^+=\Gamma\cap\Isom^+(\HH^4)\).  The orientation character shows
that \([\Gamma:\Gamma^+]=2\).  The group \(\Gamma^+\) is finitely generated
and linear in characteristic zero.  By Selberg's lemma, it contains a
torsion-free subgroup \(\Lambda\) of finite index.

The action of \(\Lambda\) on \(\HH^4\) is free and orientation preserving,
so
\[
 M=\Lambda\backslash\HH^4
\]
is an orientable hyperbolic \(4\)-manifold.  Since \(\Lambda\) has finite
index in the lattice \(\Gamma\), its quotient has finite volume.  The base
orbifold is noncompact because the census row has at least one ideal vertex;
the finite cover \(M\) is therefore noncompact as explained in
Subsection~\ref{subsec:orbifold-manifold}.

Finite-index invariance of the adjoint trace field gives
\[
 k(\Lambda)=k(\Gamma)=K_j.
\]
Theorem~\ref{thm:main} says that this field is not totally real.  Applying
Proposition~\ref{prop:total-real-obstruction} to \(\Lambda\) shows that
\(\Lambda\), and equivalently the manifold with this fundamental group, is
not pseudo-arithmetic.
\end{proof}

\section{Exact computational certification}
\label{sec:certification}

The preceding proof is hand-checkable after accepting the source rows and
the two external field theorems.  The repository also contains redundant
exact checks.  Their purpose is to detect transcription mistakes and to
make every concrete assertion reproducible.

\subsection{Rank and inertia of all three Gram matrices}

Exact symbolic elimination gives
\[
 \operatorname{rank}(G_j)=5,
 \qquad \operatorname{inertia}(G_j)=(4,1,2)
\]
for \(j=6,7,8\).  Here the entries count positive, negative, and zero
eigenvalues.  The common principal block on facets \(0,1,2,3,4\) has leading
principal minors
\begin{equation}
 1,quad 1,quad \frac{5-s}{8},quad
 \frac{3(5-s)}{32},quad -\frac{1+s}{32}.
\label{eq:leading-minors}
\end{equation}
They are all nonzero because \(1<s<3\).  The signs are
\(+,+,+,+,-\); the Jacobi signature criterion therefore gives four positive
and one negative direction on this block.  Exact elimination shows that the
full matrix has rank five, so the remaining two directions are null.  The
programs additionally verify that the determinant and every principal
six-by-six minor vanish.  The nonzero five-by-five minor and the elimination
certificate are the rank witnesses; no floating-point eigenvalue is used.

\subsection{Direct finite-volume check for the third polytope}

For \(P_{11,8}\), the script enumerates all \(127\) nonempty principal
submatrices and matches them with the ten maximal vertex subdiagrams of the
published combinatorial type.  Eight vertices have elliptic rank-four links.
The remaining two are parabolic:
\[
 \{1,3,4,6\}\quad\text{of type }\widetilde C_3,
 \qquad
 \{2,3,4,5,6\}\quad\text{of type }
 \widetilde A_1\sqcup\widetilde C_2.
\]
The second has nullity two and is the Euclidean triangular-prism link of a
nonsimple ideal vertex.  Every elliptic rank-three edge diagram has exactly
two ordinary or ideal endpoints.  Vinberg's finite-volume criterion
\cite{VinbergReflections} then gives finite volume and exactly two cusps.

An independent CoxIter run, pinned to a recorded source commit, returns
\[
 (f_0,f_1,f_2,f_3,f_4)=(10,21,18,7,1),
 \qquad \vol(P_{11,8})=\frac{19\pi^2}{864},
\]
agreeing with the census \cite{Guglielmetti}.  This redundancy verifies the
geometric input for one case; the proof of the family still cites the
Ma--Zheng classification for all three.

\subsection{A complete trace-field check for the third polytope}

The two cycles used in the theorem give only the subfield needed for the
obstruction.  Two further cycles recover \(r=\sqrt2\) and \(b_8\):
\[
 c_{126}=-r(1+s),\qquad c_{03416}=-4b_8.
\]
Since every Gram entry lies in \(\Q(s,r,a,b_8)\), this proves
\[
 k(P_{11,8})=\Q(s,r,a,b_8).
\]
With \(w=2(r+a+b_8)\), exact power-basis reconstruction gives
\(k(P_{11,8})=\Q(w)\), where
\begin{equation}
\begin{aligned}
 p_w(X)={}&X^8-8X^7-252X^6-1392X^5-3524X^4\\
          &-5280X^3-5712X^2-3968X-1136.
\end{aligned}
\label{eq:primitive-polynomial}
\end{equation}
PARI/GP verifies irreducibility, signature
\((r_1,r_2)=(4,2)\), and discriminant
\[
 40960000=2^{16}5^4.
\]
Thus the full field independently exhibits two conjugate pairs of complex
embeddings.  A second primitive-element route produces a different polynomial
for the same explicitly reconstructed field; agreement is not inferred from
degree, signature, and discriminant alone.

\subsection{Artifacts and bounded-memory execution}

The exact programs and durable outputs are:
\begin{itemize}[leftmargin=2em]
\item \path{scripts/p11_6_7_non_total_trace.py};\\
  \path{results/p11_6_7_non_total_trace_certificate.json};
\item \path{scripts/p11_8_non_total_trace.py};\\
  \path{results/p11_8_non_total_trace_certificate.json};
\item \path{scripts/p118_trace_field.py};\\
  \path{results/p118_trace_field_certificate.json};
\item \path{scripts/p118_field_independent.gp}, the independent PARI/GP
  irreducibility, real-root, signature, and discriminant check;
\item \path{scripts/p11_8_finite_volume.py};\\
  \path{results/p11_8_finite_volume_validation.json};
\item the pinned CoxIter input, output, and provenance files under
  \code{results/}.
\end{itemize}
Each expensive symbolic case runs in a fresh process.  On the reference
machine the case-worker resident-set peaks are below \(100\) MB.  The exact
commands, source URLs, versions, and expected digests are in
\code{REPRODUCIBILITY.md}.

The Ma--Zheng source file used by the certifiers has SHA-256
\[
\text{\code{0017931637f017641d4935450cc2a07d070866d9e3ede9f3a39f003}}.
\]
Each headline Python certifier checks this digest before parsing the rows.

\section{Lean formalization and trust boundary}
\label{sec:lean}

The directory \code{formal/} is a Lean~4 project pinned to Lean and mathlib
\code{v4.32.0}.  This section states exactly what its successful build means.

\subsection{Unconditional kernel-checked mathematics}

The module \code{PseudoArithmetic/Manuscript.lean} proves the following.

\begin{enumerate}[leftmargin=2.2em]
\item \code{recover\_s\_from\_twoCycle} and
  \code{recover\_u\_from\_threeCycle} check the two rearrangements in
  \eqref{eq:recover-s} and \eqref{eq:recover-u} over every
  characteristic-zero field.
\item \code{explicit\_radical\_cycle\_relation} defines the positive nested
  radicals \(s,t,a,u\) over \(\R\), proves their square and positivity facts,
  and derives \eqref{eq:cycle-relation}.
\item \code{conjugateQuadratic\_no\_real\_root} proves directly that the
  conjugate relation has no real solution when \(s^2=5\) and \(s<0\).
\item \code{not\_totallyReal\_of\_cycle\_minpoly} uses mathlib's number-field
  embedding theorem to obtain an embedding sending \(s\) to \(-\sqrt5\),
  proves that embedding nonreal, and concludes that the field is not totally
  real.
\item A \code{TwoCycleCertificate} packages exactly the four algebraic values
  and two equations used by the paper.  Its theorem \code{not\_totallyReal}
  composes the preceding results.
\end{enumerate}

The minpoly hypothesis in \code{TwoCycleCertificate} is deliberately
intrinsic: it says that the element recovered in the trace field really has
minimal polynomial \(X^2-5\).  The exact source certifier supplies this
case-specific fact; Lean then derives the bad embedding rather than accepting
that embedding as an input.

\subsection{Faithful translation of external inputs}

The structure \code{CoxeterExternalInputs} has fields for:
\begin{itemize}[leftmargin=2em]
\item the case name \(P_{11,6}\), \(P_{11,7}\), or \(P_{11,8}\);
\item the assertion that its reflection group is a finite-volume lattice;
\item the assertion that \(K\) is its adjoint trace field; and
\item the Emery--Mila implication from pseudo-arithmeticity to total reality.
\end{itemize}
The theorem \code{coxeter\_case\_main\_theorem} proves, from these fields and
a \code{TwoCycleCertificate}, the conjunction
\[
 \text{finite-volume lattice}\ \wedge\
 \text{trace-field identification}\ \wedge\
 \neg\text{total reality}\ \wedge\
 \neg\text{pseudo-arithmeticity}.
\]
Because the imported facts are structure fields supplied locally, the Lean
environment gains no axiom from this translation.

Likewise, \code{ManifoldExternalInputs} names the finite-index relation,
torsion-free orientation-preserving subgroup, quotient construction,
finite-cover geometry, trace-field invariance, and the correspondence between
a quotient and its fundamental group.  The theorem
\begin{center}
\code{finiteIndex\_manifold\_corollary}
\end{center}
checks the complete logical passage to Corollary~\ref{cor:manifolds}.

\subsection{Other formal modules}

The module \code{NonTotalReality.lean} contains the core field theorem in a
minimal form.  The separate module \code{ConditionalDescent.lean} concerns
the supplementary normalized-Clifford and split-support criterion documented
in the repository notes.  It represents ABHN, local Brauer functoriality,
Hasse--Minkowski, and the pseudo-arithmetic candidate classification as
ordinary hypotheses and checks their downstream implication graph.  That
criterion is not needed for Theorem~\ref{thm:main}; the main examples are
settled by non-total reality before any Clifford invariant is considered.

\subsection{Audit statement}

A release audit searches all project Lean sources for the declarations
\code{axiom} and \code{opaque} and for the proof escapes \code{sorry},
\code{admit}, and \code{native\_decide}.  None occurs.  The
trusted computing base is therefore the Lean kernel, the pinned Lean compiler
and mathlib declarations, and the explicitly visible hypotheses at the
external boundary.  The successful command is
\[
 \code{cd formal \&\& lake build}.
\]

\section{External-input register}
\label{sec:external-register}

This register is normative: an assertion listed here is used but not newly
proved by the manuscript or its Lean project.  Each item says both what the
paper imports and how the Lean theorem exposes it.

\begin{description}[style=nextline]
\item[\textbf{E1: Census rows and finite volume.}]
Ma--Zheng's classification supplies the three vectors, positive radical
values, cusp counts, volumes, and the assertion that they are finite-volume
hyperbolic Coxeter \(4\)-polytopes \cite{MaZheng}.  Exact scripts verify the
transcription, Gram rank, and inertia; a direct vertex-link proof independently
verifies finite volume for \(P_{11,8}\).  In Lean this is the field
\code{finite\_volume} of \code{CoxeterExternalInputs}.

\item[\textbf{E2: Vinberg field identification.}]
For a finite-volume Coxeter polytope, the field generated by cyclic products
in \(2G\) is the adjoint trace field of the reflection lattice
\cite{VinbergRings}.  In Lean the concrete algebra is a
\code{TwoCycleCertificate}; the geometric identification is the separate
field \code{trace\_field}.

\item[\textbf{E3: Pseudo-arithmetic fields are totally real.}]
Emery--Mila's definition and Section~1.4 imply that a pseudo-arithmetic
lattice in dimension \(n>3\) has totally real adjoint trace field
\cite{EmeryMila}.  Lean translates this as the structure field
\begin{center}
\code{pseudoArithmetic\_totallyReal}.
\end{center}
It is an implication, so the
kernel checks its contrapositive application without proving its geometric
premise.

\item[\textbf{E4: Extension and realization of number-field embeddings.}]
The standard extension theorem and mathlib's theorem describing evaluations
of number-field embeddings by roots of a minimal polynomial are used to
realize \(s\mapsto-\sqrt5\).  The elementary extension proof is included in
Proposition~\ref{prop:embedding-extension}; Lean uses mathlib's existing
verified result rather than a project hypothesis.

\item[\textbf{E5: Selberg's lemma.}]
A finitely generated characteristic-zero linear group has a torsion-free
finite-index subgroup \cite{Selberg}.  It is not formalized from first
principles here.  The Lean field \code{finiteIndex\_cover\_exists} exposes
the precise existence statement needed after also passing to the orientation
subgroup.

\item[\textbf{E6: Finite-index trace-field invariance.}]
The adjoint trace field is a commensurability invariant
\cite{VinbergRings}.  Lean translates the used direction as
\code{traceField\_finiteIndex\_invariant}.

\item[\textbf{E7: Finite-cover quotient geometry.}]
The standard facts that a torsion-free orientation-preserving subgroup gives
an orientable manifold, finite index preserves finite volume, and a finite
cover of the noncompact quotient remains noncompact are proved at the
elementary level in Subsection~\ref{subsec:orbifold-manifold}.  Lean packages
their geometric realization as \code{quotient\_geometry}; there is presently
no project formalization of hyperbolic orbifolds themselves.
\end{description}

No item E1--E7 is hidden behind a global Lean axiom.  Replacing an input by a
future library theorem would change the constructor used to build the
corresponding structure, not the checked proof downstream.

\section{Reproducibility and release audit}
\label{sec:reproducibility}

The repository is designed so that the central conclusion can be checked at
three increasing levels of cost.

\paragraph{Level 1: read the certificate.}
Check the four entries in \eqref{eq:four-entries}, multiply the cycles,
follow Lemmas~\ref{lem:radical-relation} and
\ref{lem:non-real-subfield}, and apply Inputs E2 and E3.  No computer algebra
is required.

\paragraph{Level 2: build the formal and written proofs.}
Run
\begin{verbatim}
cd formal
lake build

cd ../manuscript
latexmk -pdf -interaction=nonstopmode -halt-on-error \
  non_pseudo_arithmetic_4_manifolds.tex
\end{verbatim}
The project pins Lean/mathlib, and the TeX source has a fixed date.

\paragraph{Level 3: repeat all exact source checks.}
Follow \code{REPRODUCIBILITY.md} to obtain the pinned Ma--Zheng source and
HCPdm incidence file, verify their SHA-256 values, and run each Python or
PARI/GP job serially.  Reference JSON certificates are committed, so the
mathematical outputs remain inspectable without network access.  Timing and
peak-memory metadata may vary and are not used as mathematical evidence.

The command
\begin{center}
\code{python3 scripts/release\_check.py}
\end{center}
performs the local
release audit: JSON parsing, Python byte-compilation, forbidden-token checks,
documentation consistency, Lean build, exact regressions, and TeX build.  The
packaging script \code{scripts/build\_release.sh} creates a deterministic
archive and checksum under \code{dist/}.  The execution record is split
between \code{PRE\_RELEASE\_CHECKLIST.md} and \code{RELEASE\_NOTES.md}.

\appendix

\section{Dependency graph of the main conclusion}
\label{app:dependency}

For quick auditing, the complete argument is the following chain.
\[
\begin{array}{c}
\text{Ma--Zheng row (E1)}\\[2pt]
\Downarrow\\[-2pt]
\text{four entries of }A=2G\\[2pt]
\Downarrow\\[-2pt]
c_{12}=(3+s)/2,\quad c_{035}=-2u\\[2pt]
\Downarrow\\[-2pt]
s,u\in k(P),\quad \minpoly_\Q(s)=X^2-5\\[2pt]
\Downarrow\\[-2pt]
(s-1)u^2+2(s-3)u+s-7=0\\[2pt]
\Downarrow\\[-2pt]
\text{the place }s\mapsto-\sqrt5\text{ forces }u\notin\R\\[2pt]
\Downarrow\\[-2pt]
k(P)\text{ is not totally real}\\[2pt]
\Downarrow\ \text{Vinberg (E2)}\\[-2pt]
k(\Gamma)\text{ is not totally real}\\[2pt]
\Downarrow\ \text{Emery--Mila (E3)}\\[-2pt]
\Gamma\text{ is not pseudo-arithmetic}.
\end{array}
\]
The manifold corollary appends E5--E7 and E6.  The radical relation and both
downstream implication chains are checked by Lean; E1--E3 and E5--E7 remain
visible external inputs.

\section{Independent full-field data for comparison}
\label{app:full-field}

The full-field computation is not part of the proof, but it gives a useful
cross-check.  The signature \((4,2)\) of the degree-eight field means that it
has four real embeddings and two conjugate pairs of nonreal embeddings,
since \(r_1+2r_2=8\).  The two-cycle proof detects at least one of those pairs
without constructing the other six embeddings.  This explains why the short
certificate is logically stronger as an audit artifact: it contains exactly
the information needed for non-total reality and no dependence on a large
primitive-element calculation.

\begin{thebibliography}{99}

\bibitem{BBKS}
M. Belolipetsky, N. Bogachev, A. Kolpakov, and L. Slavich,
``Subspace stabilisers in hyperbolic lattices,''
\emph{Journal of the Association for Mathematical Research} 4 (2026),
111--182.
\url{https://doi.org/10.56994/JAMR.004.001.003}.

\bibitem{EmeryMila}
V. Emery and O. Mila,
``Hyperbolic manifolds and pseudo-arithmeticity,''
\emph{Transactions of the American Mathematical Society, Series B} 8 (2021),
277--295.
\url{https://doi.org/10.1090/btran/48}.

\bibitem{Guglielmetti}
R. Guglielmetti,
``CoxIter---computing invariants of hyperbolic Coxeter groups,''
\emph{LMS Journal of Computation and Mathematics} 18 (2015), 754--773.
\url{https://doi.org/10.1112/S1461157015000273}.

\bibitem{GuoMaZangZheng}
L. Guo, J. Ma, Y. Zang, and F. Zheng,
``Several families of incommensurable noncompact hyperbolic Coxeter
polytopes,'' preprint, arXiv:2607.14715, 2026.
\url{https://arxiv.org/abs/2607.14715}.

\bibitem{MaZheng}
J. Ma and F. Zheng,
``Finite-volume hyperbolic Coxeter \(4\)-dimensional polytopes with \(7\)
facets,''
\emph{Geometriae Dedicata} 219 (2025), article 25.
\url{https://doi.org/10.1007/s10711-025-00986-8}.

\bibitem{Ratcliffe}
J. G. Ratcliffe,
\emph{Foundations of Hyperbolic Manifolds}, second edition,
Graduate Texts in Mathematics 149, Springer, 2006.

\bibitem{Selberg}
A. Selberg,
``On discontinuous groups in higher-dimensional symmetric spaces,''
in \emph{Contributions to Function Theory}, Tata Institute of Fundamental
Research, Bombay, 1960, 147--164.

\bibitem{VinbergRings}
E. B. Vinberg,
``Rings of definition of dense subgroups of semisimple linear groups,''
\emph{Mathematics of the USSR-Izvestiya} 5 (1971), 45--55.
\url{https://doi.org/10.1070/IM1971v005n01ABEH001006}.

\bibitem{VinbergReflections}
E. B. Vinberg,
``Hyperbolic reflection groups,''
\emph{Russian Mathematical Surveys} 40 (1985), no.~1, 31--75.
\url{https://doi.org/10.1070/RM1985v040n01ABEH003527}.

\end{thebibliography}

\end{document}
